\section{Introduction}
Central banks change policy rates at scheduled meetings. An interest-rate model should therefore distinguish a loan that ends before a meeting from one that extends beyond it. A forward curve can express this distinction by having a different level on either side of the meeting date. As news arrives, those levels change. At the meeting itself, the announcement may cause another change.

The difficulty is to make these movements compatible with bond prices. Consider the simplest proposed model: the forward curve is constant between consecutive meetings, and each of its levels fluctuates randomly. This is easy to describe and to fit. Yet the usual no-arbitrage drift restriction does not preserve it. Even if volatility is constant across maturities within an interval, the required drift increases linearly across that interval. The curve acquires a slope. A model that insists on keeping every interval flat must eliminate that volatility.

This example illustrates the question studied here: \emph{which models of a forward curve with scheduled dates can retain their chosen shape as they evolve?} Retaining the shape is useful because it allows the entire curve to be represented by a small number of parameters. Compatibility with the drift restriction is necessary because the same evolving curve determines the prices of bonds of every maturity. We call a family \emph{consistent} when its dynamics satisfy the drift and announcement restrictions while remaining within the family.

There are two distinct ways in which a meeting enters this question. First, the curve observed today may have a break at the meeting's maturity date: rates applying before and after the meeting differ. Second, when the meeting occurs, new information may change rates for a fixed future maturity. A break in today's curve does not by itself describe an unexpected announcement. We treat these two mechanisms separately.

\subsection{Questions and results}
\paragraph{What shape can an announcement produce?}
Suppose an announcement changes the forward curve by a random level together with a maturity-dependent correction known immediately before the announcement. The correction is needed because bond prices depend exponentially on the integral of forward rates: a mean-zero rate change need not produce a mean-zero bond-price change. We show that the distribution of the random level determines the correction. If the correction is a straight line in maturity, that distribution must be Gaussian. Across several meeting intervals, the same conclusion applies after reweighting outcomes by the bond-price change accumulated over the preceding intervals. We state this reweighting explicitly in Section~\ref{sec:jumps}. It permits announcement vectors that are not jointly Gaussian. A nondegenerate hike/no-hike outcome, in contrast, is exactly incompatible with the prescribed piecewise-linear jump shape. That is a classification result rather than a large pricing effect: for the symmetric 25-basis-point example, the exact correction and its Gaussian tangent differ by only about $8\times10^{-9}$ basis points at one year. A discrete policy outcome can therefore be economically natural and numerically almost indistinguishable from the convenient affine approximation.

\paragraph{When can a fixed number of state variables describe the curve?}
Now consider continuous changes in forward rates between announcements. A common specification assigns a volatility loading according to how many meetings lie between today and the maturity. Thus one sequence of numbers specifies the sensitivity before the next meeting, between the next two meetings, and so on. A smooth maturity function, such as exponential decay, multiplies those loadings.

For any fixed finite calendar, one can keep adding state variables until the model records all past shocks that affect future maturities. A different question is whether the same state architecture remains available as the calendar grows. For one Brownian factor, constant volatility scale, and a nonzero maturity function with a finite matrix-exponential representation, we prove that it does precisely when the loading sequence satisfies a finite linear recurrence. A recurrence allows subsequent loadings to be obtained from a fixed number of preceding ones. It also allows past shocks to be summarized by a fixed number of state variables.

For example, loadings $1-\theta^i$, where $i$ is the number of intervening meetings and $0<\theta<1$, satisfy a second-order recurrence. They admit a state representation whose size does not grow with the calendar. The cumulative loadings $1+1/2+\cdots+1/i$ satisfy no finite recurrence and admit no such exact representation. Section~\ref{sec:realizations} proves both the obstruction and the construction. A complementary two-factor example in Section~\ref{sec:varianceexample} keeps only three exponential coefficients while making stochastic variance recoverable from the curve itself. Its discounted bonds are true martingales; the example shows that a stochastic coefficient on a drift-generated maturity shape need not force an unlimited enlargement of the family.

\paragraph{Can a small model still approximate a sequence with no recurrence?}
Exact representation across every calendar is a strong requirement. For a specified finite calendar, a short sum of geometric sequences may approximate the desired loadings accurately. Each such sum satisfies a recurrence. We show how its maximum loading error controls mean-square errors in forward rates and bond prices. Each model uses the drift required by its own volatility. The constants depend on the horizon and the sizes of the loadings and maturity function, without a separate dependence on the number or spacing of meetings. Section~\ref{sec:examples} applies the result to the cumulative loadings above and compares the bounds with computed errors.

\paragraph{How can meeting effects be combined with a smooth curve model?}
A model may add a curve that is linear between meetings to a smooth component, such as a sum of polynomial and exponential functions of time to maturity. The two components may share shocks. Their covariance then contributes to the required drift, so two separately admissible specifications need not remain admissible when combined.

We characterize the allowed covariance. At each maturity breakpoint, the change in the total meeting-related volatility must have zero instantaneous covariance with every nonconstant coordinate of the smooth component. A constant level is the exception. There is also a restriction within the smooth component: its covariance must be evaluated using the combined level noise of both components. With these restrictions, we give the drifts of the piecewise-linear component that satisfy the full drift equation. Section~\ref{sec:splice} develops this criterion, starting from the separate drift contributions and showing where the covariance enters. The Svensson specialization identifies the extra level covariance that the meeting component permits. For an uncorrelated block with strictly positive but potentially moving decay parameters, we also prove that its full standalone consistency equation survives: adding the meeting component does not relax the classical restrictions on those decays.

\subsection{Relation to the literature}
\paragraph{Policy meetings and overnight rates.}
The economic role of policy meetings in the term structure predates the transition to SOFR. \citet{piazzesi2005bond} models the policy target as a jump process with meeting-day effects on jump intensity. Those jumps occur at random times with elevated intensity on meeting days; this differs from an atom at an exact announcement time. \citet{das2002surprise} also allows meeting days to affect jump probabilities. Models with genuinely scheduled random-size announcements include \citet{kim2014jumps,backwell2021expected}; discrete scheduled outcomes are treated by \citet{silva2024discretely}. These models specify rate dynamics and derive prices. Our question instead fixes a maturity family and asks which dynamics preserve it.

For overnight-rate curves, \citet{gellert2021short} uses step volatilities at meeting maturities and combines the meeting component with an independent smooth component. \citet{brace2024sofr} adds stochastic volatility, exponential decay, and calibration to futures and options. Its continuous fixed-maturity forward dynamics can produce scheduled short-rate jumps by crossing maturity breakpoints. That mechanism is distinct from the additional announcement variables $\xi_n$ studied here. Our covariance criterion extends the independent combination by identifying the correlations compatible with its prescribed maturity shapes. Our finite-calendar example in Section~\ref{sec:combined} includes both diffusion and nonzero announcement changes.

\paragraph{Drift restrictions and announcement transforms.}
The continuous drift restriction is the HJM condition of \citet{heath1992bond}. \citet{fontana2018term} extends term-structure modeling to fixed discontinuity dates in a multiple-curve setting. \citet[Theorem 3.5, Remark 3.9]{fontana2022term} gives the overnight-rate formulation specialized in Section~\ref{sec:setting}. \citet{kellerressel2018affine} gives related conditions for affine semimartingales with a deterministic clock that may have atoms.

Exponential transforms and Gaussian scheduled jumps are established ingredients: see \citet{kellerressel2018affine}, \citet[Example 3.10 and Section 4.3]{fontana2022term}, and \citet[Proposition 1]{kim2014jumps}. The log-Laplace identity in Proposition~\ref{prop:profile} is a direct consequence of the jump condition. The particular conclusion developed here is the converse for a piecewise-affine jump family: each interval level must be centered Gaussian under the preceding intervals' bond-price tilt. The non-jointly-Gaussian example shows why this conclusion is weaker than requiring one jointly Gaussian announcement vector. The quantitative two-point example explains why this exact restriction can have a negligible effect at realistic rate scales.

\paragraph{Consistency of smooth curve families.}
\citet{bjork1999interest} characterizes invariance of parametric curve families, and \citet{filipovic1999note} establishes the Nelson--Siegel obstruction. \citet{filipovic2000exponential,sharef2004conditions} study exponential-polynomial families; \citet{filipovic2001consistency} develops the general consistency framework. The geometric approach to finite-dimensional invariant manifolds is developed by \citet{filipovic2004geometry}, including external stochastic factors, and extended to Wiener--Poisson jump-diffusions by \citet{filipovic2012invariant}. Scheduled maturity breakpoints require the drift identity to hold for the different analytic expressions on adjacent intervals.

The meeting component can absorb a discrepancy that is affine in maturity, so the smooth-block residual need only be affine rather than zero. With correlation it must be computed at the combined level covariance. Section~\ref{app:blocks} makes the consequence concrete: a Nelson--Siegel block can carry a stochastic level when its linear drift residual is absorbed by the front end, while its exponential shape coordinates remain restricted. The Svensson corollary likewise retains the classical exponent resonance while allowing a level correlated with its exponential-level coordinate. Theorem~\ref{thm:varyingexponents} treats a different extension, in which the decay parameters are coordinates of the state. It transfers the restrictions of \citet{filipovic2000exponential} to an uncorrelated scheduled splice with positive decays. The transfer and the explicit added covariance are the contributions; the block-alone restrictions are established results.

\paragraph{Hankel rank and realization theory.}
The linear-algebra ingredients are classical. \citet[Propositions 1--2 and Theorem 1]{hokalman1966effective} characterizes matrix sequences that admit a finite linear recurrence, identifies minimal realization dimension with Hankel rank, and constructs a realization from the sequence. Our representation of meeting loadings by powers of a matrix uses this theory.

The closest interest-rate precedent is \citet{bjork1999minimal}. For deterministic volatility that is a smooth function of time to maturity alone, its Proposition~3.1 characterizes finite-dimensional linear realizations by a matrix-exponential representation of the volatility. Proposition~4.1 identifies the minimal dimension with the span of its derivatives; Proposition~5.1 interprets the state through benchmark forward rates. The initial curve and deterministic HJM drift are supplied separately as a time-dependent output correction. Its point-process extension assumes a compensator proportional to calendar time (Assumption~7.1), so it does not include probability atoms at scheduled announcement dates.

Our volatility also depends on the number of intervening meetings. Its maturity breakpoints move as time passes, so it is outside that smooth, time-independent specification. We ask for one state dimension, one locally Lipschitz curve map, and one meeting-update rule serving every finite calendar and horizon. The additional result is that this requirement forces a recurrence of the meeting loadings, even when the curve map may be nonlinear. The sufficient construction combines continuous maturity translation with discrete meeting updates and records the HJM drift in the state. This is an extension of classical realization methods to the stated calendar convention; the recurrence algebra and the matrix-exponential maturity representation are established ingredients. Our displayed state is an upper bound, not a minimal-dimension formula. Section~\ref{sec:realizations} explains why its deterministic calendar memory must be distinguished from its random coordinates.

For smooth maturity dependence without meetings, \citet[Proposition 5.1 and Corollary 5.1]{bjork2001existence} and \citet[Corollary 8.1]{tappe2010alternative} identify the role of quasi-exponential functions. \citet{bjork2004finite} treats external stochastic-volatility states. Our sufficient construction combines those continuous-time translations with the discrete shifts of a recurrent meeting sequence. The approximation estimates use Gaussian stability and regenerate the HJM drift after replacing the loadings. They complement \citet{benth2018approximation}, whose finite-dimensional commodity-forward approximations use a Riesz basis and include convergence rates under additional regularity.

Section~\ref{sec:setting} defines the model and consistency conditions. Sections~\ref{sec:jumps}--\ref{sec:splice} develop the announcement, realization, approximation, and covariance results. Section~\ref{sec:examples} illustrates their numerical meaning. Supporting arguments are in the appendix.
